\documentclass[11pt, a4paper]{article}

% --- PACKAGES ---------------------------------------------------
\usepackage[utf8]{inputenc}
\usepackage[T1]{fontenc}
\usepackage{lmodern}
\usepackage[margin=1in]{geometry}
\usepackage{amsmath}
\usepackage{booktabs}
\usepackage{xcolor}
\usepackage{microtype}
\emergencystretch=1em
\usepackage{enumitem}
\setlist{itemsep=2pt, topsep=4pt}
\usepackage{tcolorbox}
\tcbuselibrary{skins, breakable, listings}
\usepackage{hyperref} % keep last

% --- HYPERREF ---------------------------------------------------
\hypersetup{
    colorlinks=true,
    linkcolor=blue!70!black,
    urlcolor=red!70!black,
    pdftitle={Mastering Overleaf and LaTeX: Comprehensive Study Guide},
    bookmarksopen=true
}

% --- BOXES ------------------------------------------------------
\newtcolorbox{conceptbox}[1]{
    enhanced, breakable,
    colback=blue!5!white, colframe=blue!80!black,
    title={\textbf{#1}},
    attach boxed title to top left={xshift=5mm, yshift=-3mm},
    boxed title style={colback=blue!80!black, arc=4pt},
    top=4mm, bottom=2mm, left=2mm, right=2mm, arc=4pt
}

% Listing box: content is shown verbatim, so no escaping is needed.
\newtcblisting{codebox}[1][Code Example]{
    enhanced, breakable,
    colback=gray!5!white, colframe=gray!60!black,
    title={\textbf{#1}},
    attach boxed title to top left={xshift=5mm, yshift=-3mm},
    boxed title style={colback=gray!60!black, arc=4pt},
    top=4mm, bottom=2mm, left=2mm, right=2mm, arc=4pt,
    listing only,
    listing options={basicstyle=\ttfamily\small, columns=fullflexible,
                     keepspaces=true, breaklines=true}
}

\newtcolorbox{probox}[1][Pro Tip]{
    enhanced, breakable,
    colback=yellow!10!white, colframe=orange!80!black,
    title={\textbf{#1}},
    attach boxed title to top center={yshift=-3mm},
    boxed title style={colback=orange!80!black, arc=4pt},
    top=4mm, bottom=2mm, left=2mm, right=2mm, arc=4pt
}

\newtcolorbox{warnbox}[1][Common Mistake]{
    enhanced, breakable,
    colback=red!5!white, colframe=red!70!black,
    title={\textbf{#1}},
    attach boxed title to top left={xshift=5mm, yshift=-3mm},
    boxed title style={colback=red!70!black, arc=4pt},
    top=4mm, bottom=2mm, left=2mm, right=2mm, arc=4pt
}

\newcommand{\doclink}[2]{%
    \vspace{0.2cm}\noindent
    \textbf{\textcolor{blue}{Official Documentation:}}\ \href{#1}{#2}\par}

% ================================================================
\begin{document}
% ================================================================

\begin{center}
    {\Huge \bfseries Mastering Overleaf \& \LaTeX{}}\\[0.2cm]
    {\Large \color{gray} The Comprehensive Attendee Study Guide}\\[0.5cm]
    \rule{\textwidth}{1pt}
\end{center}

Welcome to the ultimate guide to \LaTeX! This document is your self-paced
reference after the masterclass. Whether you are writing a research paper,
formatting a thesis, or building complex diagrams, this guide covers the core
concepts and the modern ``cheat codes'' to get things done effortlessly.

\vspace{0.2cm}
\tableofcontents
\vspace{0.5cm}
\noindent\rule{\textwidth}{1pt}
\newpage

% ----------------------------------------------------------------
\section{What is \LaTeX{} and Overleaf?}
% ----------------------------------------------------------------
\begin{conceptbox}{The Basics}
\begin{itemize}
    \item \textbf{\LaTeX{} (pronounced ``Lay-tech'' or ``Lah-tech'')} is a markup
    language for document preparation. Unlike Microsoft Word, where you drag
    images and fight with formatting, \LaTeX{} separates \textit{content} from
    \textit{design}. You write text with structural commands, and \LaTeX{}
    compiles it into a cleanly formatted PDF.
    \item \textbf{Overleaf} is the modern way to use \LaTeX{}. Think of it as
    Google Docs for code. It runs entirely in your browser, compiles your PDF
    next to your source, and supports real-time collaboration.
\end{itemize}
\end{conceptbox}

\doclink{https://www.overleaf.com/learn/latex/Learn_LaTeX_in_30_minutes}{Learn LaTeX in 30 Minutes}

% ----------------------------------------------------------------
\section{Document Structure}
% ----------------------------------------------------------------
Every \LaTeX{} document is split into two parts:
\begin{enumerate}
    \item \textbf{The Preamble:} The setup area at the top. Here you define the
    document class (like \texttt{article}, \texttt{report} or \texttt{book})
    and load packages (add-ons such as \texttt{graphicx} for images).
    \item \textbf{The Body:} The content, between \verb|\begin{document}| and
    \verb|\end{document}|.
\end{enumerate}

\begin{codebox}
\documentclass{article}
\usepackage{graphicx}   % <- preamble

\begin{document}        % <- body starts here

\section{Introduction}
Hello world! This is my first \textbf{bold} statement.

\end{document}
\end{codebox}

\begin{conceptbox}{Common document classes}
\begin{itemize}
    \item \texttt{article}: short papers, assignments, reports without chapters.
    \item \texttt{report}: longer documents with chapters.
    \item \texttt{book}: books, with front/back matter and two-sided layout.
    \item \texttt{beamer}: slide presentations (see Section~\ref{sec:beamer}).
    \item \texttt{IEEEtran}, \texttt{acmart}, \texttt{llncs}: publisher templates.
\end{itemize}
\end{conceptbox}

\doclink{https://www.overleaf.com/learn/latex/Creating_a_document_in_LaTeX}{Creating a document in LaTeX}

% ----------------------------------------------------------------
\section{The ``Cheat Codes'': AI and Generators}
% ----------------------------------------------------------------
\textbf{The most important lesson of this guide:} nobody memorizes \LaTeX{}
syntax anymore. Professionals use tools to generate code instantly.

\begin{conceptbox}{Modern Workflow Tools}
\begin{itemize}
    \item \textbf{ChatGPT / Claude / Gemini:} If you get a compilation error or
    need a template, ask an AI. \textit{(Prompt example: ``Give me the \LaTeX{}
    code for a 3-column table comparing Apple and Android.'')} Always paste the
    exact error message from the log.
    \item \textbf{Mathpix Snip:} Screenshot a math equation from any PDF, image
    or website and it converts it to \LaTeX{} code.
    \item \textbf{TablesGenerator.com:} An Excel-like interface. Build your table
    visually, click ``Generate'', and copy the code.
    \item \textbf{Overleaf Visual Editor:} A Word-like view of your document,
    handy for quick edits without touching the markup.
\end{itemize}
\end{conceptbox}

\begin{probox}
AI can invent package names or options that do not exist. If a generated
snippet fails, check the package documentation (\texttt{texdoc <package>} or
CTAN) before debugging blindly.
\end{probox}

\doclink{https://www.overleaf.com/learn/latex/Tables}{Tables in LaTeX}
\doclink{https://www.overleaf.com/learn/latex/Mathematical_expressions}{Mathematical expressions}

% ----------------------------------------------------------------
\section{Text Formatting, Lists and Structure}
% ----------------------------------------------------------------
\begin{conceptbox}{Everyday commands}
\begin{itemize}
    \item \verb|\textbf{bold}|, \verb|\textit{italic}|, \verb|\underline{text}|,
    \verb|\texttt{monospace}|, \verb|\emph{emphasis}|.
    \item Sizes: \verb|\tiny|, \verb|\small|, \verb|\large|, \verb|\Large|,
    \verb|\Huge|.
    \item Structure: \verb|\section|, \verb|\subsection|, \verb|\subsubsection|,
    \verb|\paragraph|.
    \item A blank line starts a new paragraph; \verb|\\| forces a line break;
    \verb|\newpage| starts a new page.
    \item Special characters need escaping: \verb|\% \& \$ \# \_ \{ \}|.
    \item Quotes: use \verb|``double''| and \verb|`single'|, never the
    straight \verb|"| key.
\end{itemize}
\end{conceptbox}

\begin{codebox}[Lists]
\begin{itemize}
    \item Bullet point
\end{itemize}

\begin{enumerate}
    \item Numbered point
\end{enumerate}

\begin{description}
    \item[Term] Its definition.
\end{description}
\end{codebox}

\begin{probox}
Prefer \verb|\emph{...}| over \verb|\textit{...}| for emphasis. It adapts
automatically, so emphasis inside an italic sentence turns upright.
\end{probox}

% ----------------------------------------------------------------
\section{Images \& Cross-Referencing}
% ----------------------------------------------------------------
In Word, moving an image by 1\,mm can destroy your entire document. In
\LaTeX{}, placement is controlled by an algorithm, and you give it hints.

\begin{codebox}
\begin{figure}[htbp]
    \centering
    \includegraphics[width=0.8\textwidth]{my_photo.jpg}
    \caption{A beautiful landscape}
    \label{fig:landscape}
\end{figure}

As seen in Figure~\ref{fig:landscape}, the mountains are tall.
\end{codebox}

\begin{probox}
Notice \verb|\label| and \verb|\ref|. \LaTeX{} renumbers figures automatically
if you insert another one before it. Use a non-breaking space
(\verb|Figure~\ref{...}|) so the word and number never split across lines.
Put \verb|\label| \emph{after} \verb|\caption|, or the number will be wrong.
\end{probox}

\begin{conceptbox}{Placement specifiers}
\begin{itemize}
    \item \texttt{h}: here if possible \quad \texttt{t}: top of page
    \quad \texttt{b}: bottom \quad \texttt{p}: separate float page.
    \item \texttt{!} overrides the rules a little; \texttt{H} (package
    \texttt{float}) forces the exact spot.
    \item Use \verb|\usepackage{graphicx}| in the preamble, and prefer PDF, PNG
    or JPG files over other formats.
\end{itemize}
\end{conceptbox}

\doclink{https://www.overleaf.com/learn/latex/Inserting_Images}{Inserting Images}

% ----------------------------------------------------------------
\section{Tables}
% ----------------------------------------------------------------
For professional-looking tables, use the \texttt{booktabs} package: no vertical
lines, and three horizontal rules (\verb|\toprule|, \verb|\midrule|,
\verb|\bottomrule|).

\begin{codebox}
\begin{table}[htbp]
    \centering
    \caption{Phone comparison}
    \label{tab:phones}
    \begin{tabular}{lcc}
        \toprule
        Feature & Apple & Android \\
        \midrule
        Open source  & No  & Mostly \\
        Price range  & High & Wide \\
        \bottomrule
    \end{tabular}
\end{table}
\end{codebox}

\noindent The result:

\begin{table}[htbp]
    \centering
    \caption{Phone comparison}
    \label{tab:phones}
    \begin{tabular}{lcc}
        \toprule
        Feature & Apple & Android \\
        \midrule
        Open source & No & Mostly \\
        Price range & High & Wide \\
        \bottomrule
    \end{tabular}
\end{table}

\begin{conceptbox}{Column types}
\texttt{l} left, \texttt{c} centre, \texttt{r} right, and \texttt{p\{3cm\}} a
fixed-width column that wraps text. Use \verb|\multicolumn| to span columns
and \verb|\multirow| (package \texttt{multirow}) to span rows.
\end{conceptbox}

% ----------------------------------------------------------------
\section{Writing Mathematics}
% ----------------------------------------------------------------
\LaTeX{} is the industry standard in STEM because its mathematical typesetting
is unmatched. Always load \texttt{amsmath}.

\begin{conceptbox}{Math Environments}
\begin{itemize}
    \item \textbf{Inline math:} wrap in single dollar signs. Typing
    \verb|$E = mc^2$| gives $E = mc^2$ inside a sentence.
    \item \textbf{Display math:} \verb|\begin{equation}| gives a centred,
    numbered equation. Use \verb|equation*| for no number.
    \item \textbf{Multi-line:} use \verb|align| to line equations up at the
    \verb|&| symbol.
\end{itemize}
\end{conceptbox}

\begin{codebox}
\begin{equation}
    x = \frac{-b \pm \sqrt{b^2 - 4ac}}{2a}
    \label{eq:quad}
\end{equation}

\begin{align}
    f(x)  &= (x+1)^2 \\
          &= x^2 + 2x + 1
\end{align}
\end{codebox}

\noindent The quadratic formula, rendered:
\begin{equation}
    x = \frac{-b \pm \sqrt{b^2 - 4ac}}{2a}
    \label{eq:quad}
\end{equation}
\noindent Refer to it with \verb|\eqref{eq:quad}|, which prints (\ref{eq:quad}).

\begin{conceptbox}{Handy symbols}
\begin{itemize}
    \item Greek: \verb|\alpha \beta \gamma \Delta \Omega| gives
    $\alpha\ \beta\ \gamma\ \Delta\ \Omega$.
    \item Operators: \verb|\sum_{i=1}^{n}|, \verb|\int_0^\infty|,
    \verb|\lim_{x\to 0}|, \verb|\partial|, \verb|\nabla|.
    \item Relations: \verb|\leq \geq \neq \approx \in \subset|.
    \item Text inside math: \verb|\text{if } x > 0| (needs \texttt{amsmath}).
\end{itemize}
\end{conceptbox}

\doclink{https://www.overleaf.com/learn/latex/Mathematical_expressions}{Mathematical expressions}

% ----------------------------------------------------------------
\section{Flowcharts and Diagrams (TikZ)}
% ----------------------------------------------------------------
How do you get flowcharts into a report?
\begin{itemize}
    \item \textbf{The Easy Way:} Draw it in Lucidchart or Draw.io, export a
    high-resolution PNG (or better, a PDF/SVG converted to PDF), and insert it
    as an image.
    \item \textbf{The Pro Way (TikZ):} \texttt{TikZ} is a \LaTeX{} graphics
    package that draws diagrams with code. It is a vector format, so it stays
    crisp at any zoom. You can even place photos inside nodes and clip them
    into circles.
\end{itemize}

\begin{codebox}[Minimal TikZ flowchart]
\usepackage{tikz}                    % preamble
\usetikzlibrary{shapes, arrows.meta, positioning}

\begin{tikzpicture}[node distance=1.6cm,
    box/.style={draw, rounded corners, minimum width=2.4cm,
                minimum height=0.9cm, fill=blue!10},
    >={Stealth}]
  \node[box] (a) {Write};
  \node[box, right=of a] (b) {Compile};
  \node[box, right=of b] (c) {PDF};
  \draw[->] (a) -- (b);
  \draw[->] (b) -- (c);
\end{tikzpicture}
\end{codebox}

\begin{warnbox}[Compile-time warning]
Large TikZ drawings are the most common cause of slow compiles. For heavy
figures, use the \texttt{standalone} class to compile the figure once to a PDF
and include that PDF as an image, or use the \verb|\tikzexternalize| library.
\end{warnbox}

\doclink{https://www.overleaf.com/learn/latex/TikZ_package}{TikZ package}

% ----------------------------------------------------------------
\section{Citations and Bibliography (BibTeX)}
% ----------------------------------------------------------------
\begin{conceptbox}{The Magic of BibTeX}
If you are writing a thesis, BibTeX saves you weeks of work:
\begin{enumerate}
    \item Create a file called \texttt{references.bib} in your project.
    \item On Google Scholar, click ``Cite'', choose ``BibTeX'', and paste that
    block into your \texttt{.bib} file.
    \item In your document, type \verb|\cite{Einstein1905}|.
\end{enumerate}
\LaTeX{} places the \texttt{[1]} in the text and builds a formatted
References section at the end.
\end{conceptbox}

\begin{codebox}[references.bib and usage]
@article{Einstein1905,
  author  = {Einstein, Albert},
  title   = {On the Electrodynamics of Moving Bodies},
  journal = {Annalen der Physik},
  year    = {1905}
}

% In the .tex file, where the references should appear:
\bibliographystyle{ieeetr}   % or plain, apalike, ...
\bibliography{references}    % file name without .bib
\end{codebox}

\begin{probox}
The modern alternative is \texttt{biblatex} with the \texttt{biber} backend.
It offers more flexible styles (APA, IEEE, Chicago) and better Unicode
support. Overleaf supports both.
\end{probox}

\doclink{https://www.overleaf.com/learn/latex/Bibliography_management_in_LaTeX}{Bibliography management in LaTeX}

% ----------------------------------------------------------------
\section{Advanced Typography \& Formatting}
% ----------------------------------------------------------------
\LaTeX{} can do things that are tedious or impossible in Word:
\begin{itemize}
    \item \textbf{Micro-typography:} The \texttt{microtype} package subtly
    adjusts character spacing so justified text avoids ugly ``rivers'' of white
    space.
    \item \textbf{Image control:} Wrap text around images (\texttt{wrapfig}),
    place images side by side (\texttt{subcaption}), or anchor an image with
    the \texttt{[H]} specifier (\texttt{float}).
    \item \textbf{Niche layouts:} Drop caps (\texttt{lettrine}), margin notes
    (\texttt{marginnote}), and text shaped into diamonds (\texttt{shapepar}).
    \item \textbf{Page layout:} \texttt{geometry} for margins, \texttt{fancyhdr}
    for headers and footers, \texttt{multicol} for columns.
\end{itemize}

% ----------------------------------------------------------------
\section{Essential Packages at a Glance}
% ----------------------------------------------------------------
\begin{center}
\begin{tabular}{@{}p{0.27\textwidth}p{0.65\textwidth}@{}}
\toprule
\textbf{Package} & \textbf{What it does} \\
\midrule
\texttt{graphicx}   & Insert images with \verb|\includegraphics|. \\
\texttt{amsmath}    & Better equations: \texttt{align}, \texttt{cases}, \verb|\text|. \\
\texttt{booktabs}   & Professional table rules. \\
\texttt{geometry}   & Set page size and margins. \\
\texttt{hyperref}   & Clickable links, references and PDF bookmarks. \\
\texttt{xcolor}     & Colours for text, boxes and backgrounds. \\
\texttt{tikz}       & Vector diagrams drawn with code. \\
\texttt{listings}   & Syntax-highlighted source code. \\
\texttt{minted}     & Higher-quality code highlighting (needs Python Pygments). \\
\texttt{siunitx}    & Consistent numbers and units, e.g.\ \verb|\SI{9.8}{m/s^2}|. \\
\texttt{algorithm2e}& Pseudocode and algorithms. \\
\texttt{subcaption} & Side-by-side sub-figures. \\
\texttt{fancyhdr}   & Custom headers and footers. \\
\texttt{csquotes}   & Correct quotation marks (recommended with \texttt{biblatex}). \\
\bottomrule
\end{tabular}
\end{center}

\begin{warnbox}[Load order]
Load \texttt{hyperref} \emph{last} (or nearly last) in the preamble. Loading
it early can break other packages. Documented exceptions, such as
\texttt{cleveref}, must come after it.
\end{warnbox}

% ----------------------------------------------------------------
\section{Presentations with Beamer}\label{sec:beamer}
% ----------------------------------------------------------------
\begin{probox}[Fun Fact]
The seminar presentation you attended was not made in PowerPoint! It was
programmed entirely in \LaTeX{} using \texttt{beamer}.
\end{probox}

\begin{codebox}[Minimal Beamer deck]
\documentclass{beamer}
\usetheme{Madrid}               % try Berlin, Warsaw, metropolis...
\title{My Talk}
\author{Your Name}
\date{\today}

\begin{document}
\begin{frame}\titlepage\end{frame}

\begin{frame}{Outline}
  \tableofcontents
\end{frame}

\section{Introduction}
\begin{frame}{Why \LaTeX{}?}
  \begin{itemize}
    \item<1-> Beautiful maths
    \item<2-> Consistent design
    \item<3-> Version-control friendly
  \end{itemize}
\end{frame}
\end{document}
\end{codebox}

\noindent The \verb|<1->| overlay syntax reveals items step by step.

% ----------------------------------------------------------------
\section{Overleaf Features Worth Knowing}
% ----------------------------------------------------------------
\begin{conceptbox}{Collaboration and workflow}
\begin{itemize}
    \item \textbf{Share:} Invite collaborators by email with edit or view access,
    or share a read-only link.
    \item \textbf{Comments and track changes:} Review a draft without editing
    it directly (availability depends on your plan).
    \item \textbf{History:} Browse and restore earlier versions of your project.
    \item \textbf{Templates gallery:} Thesis, journal, CV and poster templates
    to start from.
    \item \textbf{Git and GitHub sync:} Available on paid plans, for those who
    prefer version control outside the browser.
    \item \textbf{Menu settings:} Choose the compiler (\texttt{pdfLaTeX},
    \texttt{XeLaTeX}, \texttt{LuaLaTeX}), the TeX Live version, and the main
    document file.
\end{itemize}
\end{conceptbox}

\begin{probox}
Use \texttt{XeLaTeX} or \texttt{LuaLaTeX} (with the \texttt{fontspec}
package) when you need system fonts or non-Latin scripts such as Malayalam or
Hindi. Plain \texttt{pdfLaTeX} cannot use arbitrary system fonts.
\end{probox}

\begin{codebox}[Splitting a large project]
% main.tex
\input{preamble}                 % shared settings
\begin{document}
\include{chapters/introduction}  % starts on a new page
\include{chapters/methods}
\end{document}
\end{codebox}

% ----------------------------------------------------------------
\section{Troubleshooting Common Errors}
% ----------------------------------------------------------------
Read the \textbf{Logs and output files} panel in Overleaf. Find the
\emph{first} error. Later ones are usually a consequence of it.

\begin{center}
\begin{tabular}{@{}p{0.36\textwidth}p{0.56\textwidth}@{}}
\toprule
\textbf{Message} & \textbf{Usual cause and fix} \\
\midrule
\verb|Undefined control sequence| & A typo in a command, or the package that
defines it is not loaded. \\[4pt]
\verb|Missing $ inserted| & Maths symbols (\verb|^|, \verb|_|) used outside
math mode. Wrap them in \verb|$...$|. \\[4pt]
\verb|File 'x.sty' not found| & Wrong or unavailable package name. Check the
spelling on CTAN. \\[4pt]
\verb|Runaway argument| / \verb|File ended| & An unmatched \verb|{| or a
missing \verb|\end{...}|. \\[4pt]
\verb|Misplaced alignment tab &| & A literal \verb|&| in normal text. Write
\verb|\&|. \\[4pt]
\verb|Reference ... undefined| & Compile again; the label may be new or
misspelled. \\[4pt]
\verb|Overfull \hbox| & A line is too wide. Rephrase, or allow hyphenation
with \verb|\-|. \\[4pt]
\verb|Citation ... undefined| & The \texttt{.bib} key is wrong, or the
bibliography has not been rebuilt. \\
\bottomrule
\end{tabular}
\end{center}

\begin{probox}[Debugging trick]
Comment out half the document with \verb|%| (or the \texttt{comment} package)
and recompile. If the error disappears, it was in the half you removed. Repeat
until you find the line.
\end{probox}

% ----------------------------------------------------------------
\section{Best Practices}
% ----------------------------------------------------------------
\begin{itemize}
    \item Use \textbf{semantic commands}, e.g.\ \verb|\emph| and
    \verb|\section|, instead of manual spacing and sizing.
    \item Define repeated formatting once with \verb|\newcommand|, e.g.
    \verb|\newcommand{\R}{\mathbb{R}}| (needs \texttt{amssymb}).
    \item Always use \verb|\label| and \verb|\ref| instead of typing numbers.
    \item Keep one idea per source line, with \verb|%| comments for sections, so
    collaborators and version control can follow changes.
    \item Compile often. A small error is easier to find than fifty at the end.
    \item Keep figures in a \texttt{figures/} folder and set
    \verb|\graphicspath{{figures/}}|.
    \item Do not use \verb|\\| to make vertical space between paragraphs. Use a
    blank line and \verb|\vspace| if you need more.
\end{itemize}

% ----------------------------------------------------------------
\section{Quick Cheat Sheet}
% ----------------------------------------------------------------
\begin{codebox}[One-page reference]
\documentclass[11pt]{article}          % document type
\usepackage{graphicx,amsmath,booktabs}  % packages
\title{T}\author{A}\date{\today}\maketitle

\section{Heading}      \subsection{Sub}
\textbf{b} \textit{i} \texttt{code} \emph{e}
\begin{itemize}\item a\end{itemize}
\begin{enumerate}\item a\end{enumerate}

$a^2 + b^2 = c^2$                       % inline math
\begin{equation}\int_0^1 x\,dx\end{equation}

\includegraphics[width=\textwidth]{f.png}
\label{x}  \ref{x}  \cite{key}  \footnote{Text}
\end{codebox}

\vspace{0.5cm}
\begin{center}
    \rule{0.6\textwidth}{0.4pt}\\[4pt]
    {\small\color{gray} Happy typesetting!}
\end{center}

\end{document}
